%!TEX program = xelatex
\documentclass[a4paper,12pt]{article}
\usepackage{fontspec}\defaultfontfeatures{Ligatures=TeX}
\usepackage[a4paper,vmargin={3cm,3.5cm},hmargin={3cm,3cm}]{geometry}
%-----------------------------------------------------------------------------%
\usepackage{settings}
%-----------------------------------------------------------------------------%
%%% Title %%%
\title{G6411 Recitation Notes: Time-Series Regression Examples}
\author{Jesse Chieh Chen\\\href{mailto:cc5229@columbia.edu}{cc5229@columbia.edu}}
\date{\today}
%-----------------------------------------------------------------------------%

\begin{document}

\appendix
\section{Time-Series Regression Examples}

These examples illustrate how the serial correlation of the regression score
$g_t=X_te_t$ determines the appropriate standard errors,
and how modeling the error covariance can improve estimation.
Write $Q_{xx}=\E(X_tX_t\T)$ and let $V$ denote the asymptotic variance of
$\sqrt T(\widehat\beta-\beta)$.

\subsection{AR(1) with Conditionally Heteroskedastic Errors}

\paragraph{Model.}
Consider the stationary solution of
\begin{align}
	Y_t=\alpha Y_{t-1}+ e_t,\qquad
	 e_t=z_{t-1}z_t,\qquad
	z_t\iidto\normaldistribution(0,\sigma^2),\qquad \sigma^2>0,\quad |\alpha|<1.
\end{align}

\paragraph{Properties of Error Process.}
First, we state some properties of the error process $e_t$.
The mean, variance, and autocovariance of the errors are
\begin{align}
	\E e_t &= \E z_{t-1} z_t = \E z_{t-1}\E z_t = 0, \\
	\E e_t^2 &= \E z_{t-1}^2 z_t^2 = \E z_{t-1}^2 \E z_t^2 = \sigma^4, \\
	\E e_t e_{s} &= \E z_{t-1} z_t z_{s-1} z_s = \E z_{t} \E(z_{t-1} z_{s}) \E z_{s-1} = 0 \quad\text{for}\quad t > s.
\end{align}
Hence, the error process is white noise.
Furthermore, it is strictly stationary and ergodic.
Let $\mathcal F_t=\sigma(z_t,z_{t-1},\ldots)$ be the information generated by all shocks through date $t$.
Since $z_t$ is independent of $\mathcal F_{t-1}$,
\begin{align}
	\E(e_t\given\mathcal F_{t-1})=0,
	\qquad \Var(e_t\given\mathcal F_{t-1})=\sigma^2z_{t-1}^2.
\end{align}
The conditional variance varies with the previous shock, so the errors are conditionally heteroskedastic.
Their unconditional variance is nevertheless the same $\sigma^4$ at every date:
there is no unconditional heteroskedasticity.

\paragraph{Properties of the Dependent Variable.}
First, notice that we have
\begin{align}
	\E Y_{t-1} e_t = \E Y_{t-1} z_{t-1} z_t = \E (Y_{t-1} z_{t-1}) \E z_t = 0.
\end{align}
Stationarity of $\{Y_t\}$ gives a common mean $\mu$ and variance $\gamma(0)$, hence we have
\begin{align}
	\E Y_t &= \E(\alpha Y_{t-1} + e_t) = \alpha \E Y_{t-1} + \E e_t \nonumber\\
	&\implies
	\mu = \alpha \mu + 0 \implies \mu = 0,\\
	\Var(Y_t) &= \E(\alpha Y_{t-1} + e_t)^2 = \alpha^2 \E Y_{t-1}^2 + \E e_t^2 + 2\alpha \E(Y_{t-1} e_t) \nonumber\\
	&\implies \gamma(0) = \alpha^2 \gamma(0) + \sigma^4 \implies \gamma(0) = \frac{\sigma^4}{1 - \alpha^2}.
\end{align}
where we used the fact that $\E(Y_{t-1} e_t) = 0$ and $\E e_t^2 = \sigma^4$.
We shall assume that $Y_t$ is also strictly stationary and ergodic from now on,
which is true for the AR(1) process with $|\alpha|<1$.

\paragraph{Estimation.}
Regress $Y_t$ on $Y_{t-1}$, writing $X_t=Y_{t-1}$.
The \gls{ols} estimator without an intercept is
\begin{align}
	\widehat\alpha
	&= \left(\sum_{t=1}^{T} X_t X_t\T\right)\inv \sum_{t=1}^{T} X_t Y_t \\
	&= \left(\frac1T\sum_{t=1}^{T} Y_{t-1}^2\right)\inv \frac1T\sum_{t=1}^{T} Y_{t-1}Y_t
	= \alpha + \left(\frac1T\sum_{t=1}^{T} Y_{t-1}^2\right)\inv \frac1T\sum_{t=1}^{T} Y_{t-1}e_t.
\end{align}
Rearranging gives
\begin{align}
	\sqrt{T}(\widehat\alpha - \alpha)
	= \left(\frac1T\sum_{t=1}^{T} Y_{t-1}^2\right)\inv \frac1{\sqrt{T}}\sum_{t=1}^{T} g_t
	\quad\text{where}\quad
	g_t = X_t e_t = Y_{t-1} e_t.
\end{align}

\paragraph{Properties of Score.}
We verify that the score $g_t$ forms a \gls{mds} relative to $\{\mathcal F_t\}$.
Since $Y_{t-1}$ and $z_{t-1}$ are known at date $t-1$, we have
\begin{align}
	\E(g_t\given\mathcal F_{t-1})
	&= \E(Y_{t-1} e_t\given\mathcal F_{t-1}) \notag\\
	&= \E(Y_{t-1} z_{t-1} z_t\given\mathcal F_{t-1})
	= Y_{t-1} z_{t-1} \E(z_t\given\mathcal F_{t-1}) \notag\\
	&= Y_{t-1} z_{t-1} \cdot 0 = 0.
\end{align}
Thus, $g_t$ is an \gls{mds}, despite the dependence in $Y_t$ and in $ e_t$.
Now we calculate the variance of the score $g_t$:
\begin{align}
	\Omega
	\coloneqq \E g_t^2
	&= \E(Y_{t-1}^2 e_t^2)
	= \E(Y_{t-1}^2 z_{t-1}^2 z_t^2) \notag\\
	&= \E(Y_{t-1}^2 z_{t-1}^2) \E(z_t^2)
	= \E(Y_{t-1}^2 z_{t-1}^2) \sigma^2.
\end{align}
Now focus on the term $\E(Y_{t-1}^2 z_{t-1}^2)$:
\begin{align}
	\E(Y_{t-1}^2 z_{t-1}^2)
	&= \E((\alpha Y_{t-2} + e_{t-1})^2 z_{t-1}^2) \notag\\
	&= \alpha^2\E(Y_{t-2}^2 z_{t-1}^2) + 2\alpha \E(Y_{t-2} e_{t-1} z_{t-1}^2)
	+ \E(e_{t-1}^2 z_{t-1}^2) \notag\\
	&= \alpha^2\E(Y_{t-2}^2)\E(z_{t-1}^2)
	+2\alpha\E(Y_{t-2}z_{t-2})\E(z_{t-1}^3)
	+\E(z_{t-2}^2)\E(z_{t-1}^4) \notag\\
	&= \alpha^2\frac{\sigma^4}{1-\alpha^2}\sigma^2+0+\sigma^2\cdot3\sigma^4
	=\sigma^6\left(\frac{\alpha^2}{1-\alpha^2}+3\right).
\end{align}
Here we used $e_{t-1}=z_{t-2}z_{t-1}$, independence of $z_{t-1}$ from $(Y_{t-2},z_{t-2})$,
and the normal moments $\E z_{t-1}^3=0$ and $\E z_{t-1}^4=3\sigma^4$.
Consequently,
\begin{align}
	\Omega=\sigma^8\left(\frac{\alpha^2}{1-\alpha^2}+3\right)
	=\sigma^8\frac{3-2\alpha^2}{1-\alpha^2}.
\end{align}

\paragraph{Asymptotic Distribution.}
Since $\{Y_t\}$ is strictly stationary and ergodic,
the ergodic theorem gives
\begin{align}
	\widehat Q_{xx} =
	\frac1T\sum_{t=1}^{T} Y_{t-1}^2 \pto \E Y_{t-1}^2 = Q_{xx} = \gamma(0) = \frac{\sigma^4}{1 - \alpha^2}.
\end{align}
The score $g_t$ is a function of the current and past iid shocks, so it is also strictly stationary and ergodic.
Since it is an \gls{mds} with finite variance, the \gls{mds} \gls{clt} applies and gives
\begin{align}
	\frac1{\sqrt{T}}\sum_{t=1}^{T} g_t \dto \normaldistribution(0, \Omega).
\end{align}
Combining them, we have
\begin{gather}
	\sqrt T(\widehat\alpha-\alpha)
	=\widehat Q_{xx}^{-1}\frac1{\sqrt T}\sum_{t=1}^T g_t
	\dto Q_{xx}\inv\normaldistribution\bigl(0,\Omega\bigr)
	= \normaldistribution\bigl(0,V\bigr)
	\shortintertext{where}
	V
	= Q_{xx}\inv \Omega Q_{xx}\inv
	= \left(\frac{\sigma^4}{1-\alpha^2}\right)^{-2} \sigma^8 \left(\frac{\alpha^2}{1-\alpha^2}+3\right)
	= (1-\alpha^2)(3-2\alpha^2).
\end{gather}

\paragraph{Inference.}
Define $\widehat e_t=Y_t-\widehat\alpha Y_{t-1}$ and $\widehat g_t=Y_{t-1}\widehat e_t$.
The HC0 estimator
\begin{gather}
	\widehat V = \widehat Q_{xx}\inv \widehat\Omega_{\mathrm{HC0}} \widehat Q_{xx}\inv
	\quad\text{where}\quad
	\widehat\Omega_{\mathrm{HC0}}=\frac1T\sum_{t=1}^T \widehat{g}_t \widehat{g}_t\T
\end{gather}
is a consistent estimator of $V$.
Conditional heteroskedasticity calls for the sandwich variance,
but no lagged score products are needed: HC0 suffices because $g_t$ is an \gls{mds}.

\subsection{Exogenous Regressors, but Serially Correlated Scores}

\paragraph{Model.}
Consider the stationary solutions of
\begin{align}
	Y_t=\beta_0+\beta_1Z_t+e_t,\qquad
	Z_t=\delta Z_{t-1}+v_t,\qquad e_t=\rho e_{t-1}+u_t,
\end{align}
where $0<\delta,\rho<1$.
Let $\{v_t\}$ and $\{u_t\}$ be independent sequences of independent normal shocks,
with mean zero and variances $1-\delta^2$ and $1-\rho^2$, respectively.

\paragraph{Properties of Error Process.}
Stationarity and independence of $u_t$ from past errors give
\begin{align}
	\E e_t&=\rho\E e_{t-1}+\E u_t
	\quad\implies\quad \E e_t=0,\\
	\Var(e_t)&=\rho^2\Var(e_{t-1})+1-\rho^2
	\quad\implies\quad \Var(e_t)=1.
\end{align}
For $j\ge1$, the same independence gives the autocovariances recursively:
\begin{align}
	\E(e_te_{t-j})
	=\rho\E(e_{t-1}e_{t-j})+\E(u_te_{t-j})
	=\rho\E(e_{t-1}e_{t-j})=\rho^j.
\end{align}
Thus the errors are serially correlated, although their unconditional variance is constant.
Let $\mathcal F_t=\sigma(u_s,v_s:s\le t)$ be the information generated by both shock histories through date $t$.
Then
\begin{align}
	\E(e_t\given\mathcal F_{t-1})=\rho e_{t-1},
	\qquad \Var(e_t\given\mathcal F_{t-1})=1-\rho^2.
\end{align}
The errors are conditionally homoskedastic, but they are not an \gls{mds}.
The stationary error process is strictly stationary and ergodic.

\paragraph{Properties of the Regressor and Dependent Variable.}
The same calculations for $Z_t$ give
\begin{align}
	\E Z_t=0,\qquad \Var(Z_t)=1,\qquad \E(Z_tZ_{t-j})=\delta^j,\quad j\ge0.
\end{align}
Since the two shock sequences are independent, the entire regressor process is independent of the entire error process.
In particular,
\begin{align}
	\E(e_t\given\{Z_s\}_{s\in\mathbb Z})=0,
	\qquad \E(Z_te_t)=\E Z_t\E e_t=0.
\end{align}
It also follows that
\begin{align}
	\E Y_t&=\beta_0+\beta_1\E Z_t+\E e_t=\beta_0,\\
	\Var(Y_t)&=\beta_1^2\Var(Z_t)+\Var(e_t)+2\beta_1\E(Z_te_t)=\beta_1^2+1.
\end{align}
The regressor and dependent variable are jointly strictly stationary and ergodic.

\paragraph{Estimation.}
Regress $Y_t$ on an intercept and $Z_t$, writing
$X_t=(1,Z_t)\T$ and $\beta=(\beta_0,\beta_1)\T$.
The \gls{ols} estimator is
\begin{align}
	\widehat\beta
	=\left(\frac1T\sum_{t=1}^T X_tX_t\T\right)\inv\frac1T\sum_{t=1}^T X_tY_t
	=\beta+\widehat Q_{xx}\inv\frac1T\sum_{t=1}^T X_te_t.
\end{align}
Here $\widehat Q_{xx}=T^{-1}\sum_{t=1}^T X_tX_t\T$.
Rearranging gives
\begin{align}
	\sqrt T(\widehat\beta-\beta)
	=\widehat Q_{xx}\inv\frac1{\sqrt T}\sum_{t=1}^T g_t,
	\qquad g_t=X_te_t=\begin{pmatrix}e_t\\Z_te_t\end{pmatrix}.
\end{align}

\paragraph{Properties of Score.}
The score has mean zero, since $\E e_t=0$ and $\E(Z_te_t)=0$.
However, conditional on $\mathcal F_{t-1}$, the new shocks $u_t$ and $v_t$ are independent and have mean zero, so
\begin{align}
	\E(g_t\given\mathcal F_{t-1})
	=
	\begin{pmatrix}
		\rho e_{t-1}\\
		\E[(\delta Z_{t-1}+v_t)(\rho e_{t-1}+u_t)\given\mathcal F_{t-1}]
	\end{pmatrix}
	=
	\begin{pmatrix}
		\rho e_{t-1}\\
		\delta\rho Z_{t-1}e_{t-1}
	\end{pmatrix}.
\end{align}
Thus $g_t$ is not an \gls{mds}.
Nevertheless, its autocovariances have a simple form:
\begin{align}
	\Gamma_g(j)
	\coloneqq\E(g_tg_{t-j}\T)
	&=\E(X_tX_{t-j}\T)\E(e_te_{t-j})\notag\\
	&=\begin{pmatrix}1&\E Z_{t-j}\\\E Z_t&\E(Z_tZ_{t-j})\end{pmatrix}\rho^j
	=\begin{pmatrix}\rho^j&0\\0&(\delta\rho)^j\end{pmatrix},\qquad j\ge0.
\end{align}
In particular, the contemporaneous score variance is $\Gamma_g(0)=I_2$,
but the \gls{lrv} also includes the nonzero-lag covariances.
Since $\Gamma_g(-j)=\Gamma_g(j)\T$, summing the geometric series gives
\begin{align}
	\Omega
	\coloneqq\sum_{j=-\infty}^{\infty}\Gamma_g(j)
	&=I_2+2\sum_{j=1}^{\infty}\begin{pmatrix}\rho^j&0\\0&(\delta\rho)^j\end{pmatrix}\notag\\
	&=\begin{pmatrix}
		1+\dfrac{2\rho}{1-\rho}&0\\[6pt]
		0&1+\dfrac{2\delta\rho}{1-\delta\rho}
	\end{pmatrix}
	=\begin{pmatrix}
		\dfrac{1+\rho}{1-\rho}&0\\[6pt]
		0&\dfrac{1+\delta\rho}{1-\delta\rho}
	\end{pmatrix}.
\end{align}

\paragraph{Asymptotic Distribution.}
The ergodic theorem gives
\begin{align}
	\widehat Q_{xx}\pto Q_{xx}
	=\E\begin{pmatrix}1&Z_t\\Z_t&Z_t^2\end{pmatrix}=I_2.
\end{align}
It also gives $T^{-1}\sum_tg_t\pto\E g_t=0$, so $\widehat\beta\pto\beta$.
The stable Gaussian models have finite moments of every order and geometrically decaying mixing coefficients.
Their score therefore satisfies a dependent-data \gls{clt}:
\begin{align}
	\frac1{\sqrt T}\sum_{t=1}^T g_t\dto\normaldistribution(0,\Omega).
\end{align}
Combining these results,
\begin{gather}
	\sqrt T(\widehat\beta-\beta)
	\dto\normaldistribution(0,V)
	\quad\text{where}\quad
	V=Q_{xx}\inv\Omega Q_{xx}\inv=\Omega
	=\begin{pmatrix}
		\dfrac{1+\rho}{1-\rho}&0\\[6pt]
		0&\dfrac{1+\delta\rho}{1-\delta\rho}
	\end{pmatrix}.
\end{gather}

\paragraph{Inference.}
Define $\widehat e_t=Y_t-X_t\T\widehat\beta$ and $\widehat g_t=X_t\widehat e_t$.
The HC0 estimator $T^{-1}\sum_t\widehat g_t\widehat g_t\T$ estimates only $\Gamma_g(0)=I_2$,
missing the nonzero-lag contributions to $\Omega$.
It therefore understates the asymptotic variances of both coefficients.
To estimate these contributions without fitting a parametric covariance model,
use the \gls{nw} estimator with an appropriate bandwidth $L_T$:
\begin{gather}
	\widehat V_{\mathrm{NW}}=\widehat Q_{xx}\inv\widehat\Omega_{\mathrm{NW}}\widehat Q_{xx}\inv\\
	\shortintertext{where}
	\widehat\Omega_{\mathrm{NW}}
	=\widehat\Gamma_g(0)+\sum_{j=1}^{L_T}\left(1-\frac{j}{L_T+1}\right)
	[\widehat\Gamma_g(j)+\widehat\Gamma_g(j)\T],\\
	\widehat\Gamma_g(j)=\frac1T\sum_{t=j+1}^T\widehat g_t\widehat g_{t-j}\T.
\end{gather}

\subsection{An Alternative Approach: \glsentryshort{gls}}

\paragraph{Model and Motivation.}
Keep the model and assumptions of the previous example:
\begin{align}
	Y_t=X_t\T\beta+e_t,\qquad X_t=(1,Z_t)\T,\qquad
	e_t=\rho e_{t-1}+u_t.
\end{align}
The \gls{nw} estimator corrects the standard errors without changing the \gls{ols} coefficients.
If we are willing to use the AR(1) model for the errors, we can instead remove their serial correlation
before estimating $\beta$.
This leads to \gls{gls}.

\paragraph{Transforming the Regression.}
First suppose $\rho$ is known.
Subtract $\rho$ times the equation at $t-1$ from the equation at $t$:
\begin{align}
	Y_t-\rho Y_{t-1}
	&=(X_t-\rho X_{t-1})\T\beta+e_t-\rho e_{t-1}\notag\\
	&=(1-\rho)\beta_0+(Z_t-\rho Z_{t-1})\beta_1+u_t,
	\qquad t=2,\ldots,T.
\end{align}
Define $\widetilde Y_t=Y_t-\rho Y_{t-1}$ and
$\widetilde X_t=X_t-\rho X_{t-1}$.
The transformed equation is $\widetilde Y_t=\widetilde X_t\T\beta+u_t$,
with independent errors of variance $1-\rho^2$.
This transformation is called \emph{quasi-differencing}.
Notice that the intercept column becomes $1-\rho$, not one.
For exact \gls{gls}, retain the first observation by setting
\begin{align}
	\widetilde Y_1=\sqrt{1-\rho^2}\,Y_1,
	\qquad \widetilde X_1=\sqrt{1-\rho^2}\,X_1.
\end{align}
Its error is $\widetilde e_1=\sqrt{1-\rho^2}\,e_1$;
for $t\ge2$, write $\widetilde e_t=u_t$.
These transformed errors all have variance $1-\rho^2$ and are independent of one another
and of the entire transformed regressor process.
Thus their conditional covariance matrix is $(1-\rho^2)I_T$.

\paragraph{Estimation.}
Apply \gls{ols} to the transformed regression:
\begin{align}
	\widehat\beta_{\mathrm{GLS}}
	=\left(\frac1T\sum_{t=1}^T\widetilde X_t\widetilde X_t\T\right)\inv
	\frac1T\sum_{t=1}^T\widetilde X_t\widetilde Y_t
	=\beta+\widehat Q_{\widetilde x\widetilde x}\inv
	\frac1T\sum_{t=1}^T\widetilde X_t\widetilde e_t.
\end{align}
Here $\widehat Q_{\widetilde x\widetilde x}=T^{-1}\sum_{t=1}^T\widetilde X_t\widetilde X_t\T$.

\paragraph{Properties of Score.}
For $t\ge2$, the score is $\widetilde g_t=\widetilde X_tu_t$.
Although $\widetilde X_t$ contains the current shock $v_t$,
$u_t$ is independent of $(\mathcal F_{t-1},v_t)$ and has mean zero.
Therefore,
\begin{align}
	\E(\widetilde g_t\given\mathcal F_{t-1})
	&=\E\left[\widetilde X_t\E(u_t\given\mathcal F_{t-1},v_t)
	\given[\middle]\mathcal F_{t-1}\right]=0.
\end{align}
The score $\widetilde g_t$ is an \gls{mds}.
Since $u_t$ is independent of $\widetilde X_t$, the variance of the score is
\begin{gather}
	\widetilde\Omega
	\coloneqq\E(\widetilde g_t\widetilde g_t\T)
	=\E(\widetilde X_t\widetilde X_t\T u_t^2)
	=\E(\widetilde X_t\widetilde X_t\T)\E u_t^2
	=(1-\rho^2)Q_{\widetilde x\widetilde x}
	\shortintertext{where}
	Q_{\widetilde x\widetilde x}
	\coloneqq\E(\widetilde X_t\widetilde X_t\T)
	=\begin{pmatrix}(1-\rho)^2&0\\0&\E(Z_t-\rho Z_{t-1})^2\end{pmatrix}
\end{gather}
with $\E(Z_t-\rho Z_{t-1})^2=1+\rho^2-2\rho\delta$.
Hence, we have
\begin{align}
	\widetilde\Omega
	&=(1-\rho^2)\begin{pmatrix}(1-\rho)^2&0\\0&1+\rho^2-2\rho\delta\end{pmatrix}.
\end{align}

\paragraph{Asymptotic Distribution.}
The ergodic theorem gives $\widehat Q_{\widetilde x\widetilde x}\pto Q_{\widetilde x\widetilde x}$.
The first observation contributes only $o_p(1)$ to the normalized score sum,
so the \gls{mds} \gls{clt} and Slutsky's theorem give
\begin{gather}
	\sqrt T(\widehat\beta_{\mathrm{GLS}}-\beta)
	\dto\normaldistribution(0,V_{\mathrm{GLS}})
	\shortintertext{where}
	V_{\mathrm{GLS}}
	=Q_{\widetilde x\widetilde x}\inv\widetilde\Omega Q_{\widetilde x\widetilde x}\inv
	=(1-\rho^2)Q_{\widetilde x\widetilde x}\inv
	=\begin{pmatrix}
		\dfrac{1+\rho}{1-\rho}&0\\[6pt]
		0&\dfrac{1-\rho^2}{1+\rho^2-2\rho\delta}
	\end{pmatrix}.
\end{gather}
In particular, $\widehat\beta_{\mathrm{GLS}}\pto\beta$.
Compared with \gls{ols}, the intercept has the same asymptotic variance, but the slope has a smaller one:
\begin{align}
	V_{\mathrm{OLS},22}-V_{\mathrm{GLS},22}
	=\frac{1+\delta\rho}{1-\delta\rho}
	-\frac{1-\rho^2}{1+\rho^2-2\rho\delta}
	=\frac{2\rho^2(1-\delta^2)}{(1-\delta\rho)(1+\rho^2-2\rho\delta)}>0.
\end{align}
Thus, modeling the error correlation can make the estimation more efficient.

\paragraph{Inference and Unknown \texorpdfstring{$\rho$}{rho}.}
In practice, estimate $\rho$ using the residuals $\widehat e_t$ from the original \gls{ols} regression:
\begin{align}
	\widehat\rho=\frac{\sum_{t=2}^T\widehat e_t\widehat e_{t-1}}
	{\sum_{t=2}^T\widehat e_{t-1}^2}\pto\rho.
\end{align}
Once we have $\widehat\rho$, we can perform the quasi-differencing transformation using it.
This is the \gls{fgls} procedure.
Under the present stationary Gaussian model, estimating $\rho$ does not change the limiting distribution above.
Estimate the transformed error variance and the asymptotic variance by
\begin{align}
	\widehat V_{\mathrm{FGLS}}
	&=\widehat\sigma_u^2\left(\frac1T\sum_{t=1}^T
	\widehat{\widetilde X}_t\widehat{\widetilde X}_t\T\right)\inv
	\quad\text{where}\quad
	\widehat\sigma_u^2
	=\frac1{T}\sum_{t=1}^T
	\left(\widehat{\widetilde Y}_t-\widehat{\widetilde X}_t\T\widehat\beta_{\mathrm{FGLS}}\right)^2
\end{align}
where $\widehat{\widetilde Y}_t$ and $\widehat{\widetilde X}_t$ denote the resulting transformed data.

\vfill
\section*{Acronyms}
\printglossaries

\end{document}
